\scanpage{19}
\uline{Remarque} On n'a jamais, dans les calculs précédents, utilisé la structure additive des $A(X)$, seulement leur structure de monoïde. Donc il y a lieu de formuler la théorie des correspondances en prenant \deleted{\unsure{des}} $X\longrightarrow A(X)$ \deleted{qui sont des} \inserted{foncteurs dans les} ``monoïdes bivariants'' \ldots.

\deleted{Il faudrait être \deleted{\unsure{prudent}} \inserted{\unsure{rigoureux}} pour les signes, dans le cas de monoïdes non commutatifs, en prenant des monoïdes gradués [\unsure{les degrés en} \unsure{anneaux anticommutatifs}}

Ce point de vue cependant ne semble pas très commode dans le cas ou les 2 formules de projections \inserted{(gauche et droite)} ne sont pas valables l'une et l'autre, et qu'il faut introduire des signes. Il faudrait alors remplacer les monoïdes par des ``monoïdes jaugés'' $(A^i)_{i\in\mathbb Z}$, loi de composition $A^i\times\unsure{A^j}-A^{i+j}$ (NB: L'ensemble somme des $A^i$ \deleted{\unsure{est}} muni d'une structure de monoïde) et \deleted{\unsure{automorphismes involutifs d'ordre}} \inserted{applications bijectives} $\sigma^i:A^i-A^i$ satisfaisant $\unsure{\sigma^{i\,2}|(\sigma^i)^2}=\operatorname{id}$, $\sigma^i=\operatorname{id}$ si $i$ pair, \deleted{\unsure{pour}} enfin chaque $A^*$ muni d'un \uline{degré} \inserted{$d(A^*)$} qui permet d'imposer le degré des $\unsure{f_*|f_s}$ des $(\unsure{f_*,f^*|f=(f^i)}):A^*\unsure{\longrightarrow|-}B^*$ comme $\deleted{\unsure{f|d}}d(B^*)-d(A^*)$. On peut alors \unsure{indiquer} comment les \uline{signes} doivent être calculés dans la formule de projection gauche (si on veut que la droite soit \unsure{valable}).
\mspage{18}
